\begin{helper}
Let $G$ be a finite $\Delta$-regular graph, let $\mu$ be the
activation--priority independent-set sample of
\noderef{RandomIndependentSetSampling}, and fix a vertex $r$ and a set
$X\subseteq N_G(r)$ such that no two distinct vertices of $X$ have a common
neighbor outside $N_G[r]$.  Let $Q=\{a,b,c\}\subseteq X$ be an independent
triple.  For each independent pair $P=\{p,q\}\subseteq X$, let
\[
\ell_P=\frac{|N_G(p)\cap N_G(q)|}{\Delta}.
\]
Then
\[
\Pr[Q\subseteq I]
\le
\frac1{\Delta^3}
+\frac{1}{3\Delta^3}\sum_{\{p,q\}\subseteq Q}
\left(\frac{2}{(2-\ell_{\{p,q\}})(1-\ell_{\{p,q\}})}-1\right)
\]
\end{helper}
\begin{proof}
Write $Q=\{a,b,c\}$.  First fix an independent pair $\{p,q\}\subseteq Q$.
Because $Q\subseteq X$ and $X\subseteq N_G(r)$, both $p$ and $q$ are
neighbors of $r$.  Because $Q$ is independent, $p\ne q$ and $p$ is not
adjacent to $q$.  If $w$ is a common neighbor of $p$ and $q$, then either
$w=r$ or $w$ is adjacent to $r$: otherwise $w\ne r$ and $w\notin N_G(r)$
would contradict the clean-neighborhood hypothesis for the two distinct
vertices $p,q\in X$.  Also $w\ne p,q$, since the graph is simple and
$p,q$ are non-adjacent.  Hence
\[
N_G(p)\cap N_G(q)\subseteq \{r\}\cup (N_G(r)\setminus\{p,q\}).
\]
The right side has at most $\Delta-1$ vertices, since $G$ is
$\Delta$-regular and $p,q$ are two distinct neighbors of $r$.  Thus
\[
0\le \ell_{\{p,q\}}\le 1-\frac1\Delta<1,
\]
and in particular each
$\lambda_{pq}:=1-\ell_{\{p,q\}}$ is strictly positive.

Condition on the event that all three vertices $a,b,c$ are activated.  The
activation part of \noderef{RandomIndependentSetSampling} gives the factor
$(\gamma/\Delta)^3$.  Put $x_v=\gamma(1-\pi(v))$ for $v\in Q$.  The
independent uniform-priority law in \noderef{RandomIndependentSetSampling}
gives conditional density $\gamma^{-3}\,dx_a\,dx_b\,dx_c$ on the box
$[0,\gamma]^3$ for these three variables, so the activation factor and the
Jacobian leave the outside factor $\Delta^{-3}$.  Since
\noderef{RandomIndependentSetSampling} is a Bernoulli activation law with
probability $\gamma/\Delta$, we have $0\le\gamma/\Delta\le1$; hence
$0\le x_v/\Delta\le1$ throughout the box.

Now fix $x_a,x_b,x_c$.  Let $w\notin Q$ have nonempty adjacency pattern
$S\subseteq Q$.  Conditional on the three priorities in $Q$, the vertex $w$
defeats the triple exactly when it is activated and has priority larger than
at least one of the vertices of $S$.  In the variables
$x_v=\gamma(1-\pi(v))$, this event has conditional probability
\[
\frac{\gamma}{\Delta}\left(1-\min_{v\in S}\pi(v)\right)
=\frac{\max_{v\in S}x_v}{\Delta}.
\]
The activation and priority coordinates of distinct outside vertices are
independent by \noderef{RandomIndependentSetSampling}.  Therefore the
conditional probability that no outside vertex defeats the triple is the
product, over all outside vertices with nonempty pattern $S$, of
\[
1-\frac{\max_{v\in S}x_v}{\Delta}.
\]

Consider the ordered chamber $x_a\ge x_b\ge x_c$.  Classify vertices outside
$Q$ by their nonempty adjacency pattern to $Q$, writing
$n_a,n_b,n_c,n_{ab},n_{ac},n_{bc},n_{abc}$ for the seven pattern counts.
The clean-neighborhood hypothesis shows that the common neighbors counted by
$\ell_{ab}$ are precisely the vertices in the $ab$ and $abc$ classes, and
cyclically
\[
\ell_{ab}=\frac{n_{ab}+n_{abc}}{\Delta},\quad
\ell_{ac}=\frac{n_{ac}+n_{abc}}{\Delta},\quad
\ell_{bc}=\frac{n_{bc}+n_{abc}}{\Delta}.
\]
On this chamber, the classes $a,ab,ac,abc$ have maximum coordinate $x_a$,
the classes $b,bc$ have maximum coordinate $x_b$ after the classes already
charged to $x_a$ are removed, and the class $c$ has maximum coordinate
$x_c$.  Since $G$ is $\Delta$-regular and the vertices of $Q$ are pairwise
non-adjacent, the $x_a$ exponent is
$n_a+n_{ab}+n_{ac}+n_{abc}=\Delta$, the $x_b$ exponent is
$n_b+n_{bc}=\Delta-n_{ab}-n_{abc}=(1-\ell_{ab})\Delta$, and the normalized
$x_c$ exponent is
\[
\rho_c
=\frac{n_c}{\Delta}
=1-\ell_{ac}-\ell_{bc}+\frac{n_{abc}}{\Delta}.
\]
The exact product on this chamber is therefore
\[
\left(1-\frac{x_a}{\Delta}\right)^\Delta
\left(1-\frac{x_b}{\Delta}\right)^{(1-\ell_{ab})\Delta}
\left(1-\frac{x_c}{\Delta}\right)^{\rho_c\Delta}.
\]
Using $1-t\le e^{-t}$ for $0\le t\le1$, it is at most
\[
\exp(-x_a)\exp(-(1-\ell_{ab})x_b)\exp(-\rho_c x_c),
\]
Enlarging the chamber from
$[0,\gamma]^3$ to the infinite chamber only increases the integral.

The two chambers in which $a$ and $b$ are the two larger coordinates
therefore contribute at most
\[
C_{ab}=2\int_{0\le z\le y\le x<\infty}
e^{-x}e^{-(1-\ell_{ab})y}e^{-\rho_c z}\,dz\,dy\,dx .
\]
Writing $\lambda_{ab}=1-\ell_{ab}$, direct integration gives
\[
C_{ab}=\frac{2}{(1+\lambda_{ab})(1+\lambda_{ab}+\rho_c)},
\]
with the case $\rho_c=0$ interpreted by the limiting value.  The same
argument gives analogous bounds $C_{ac}$ and $C_{bc}$.

Put
\[
x=\lambda_{ab},\qquad y=\lambda_{ac},\qquad z=\lambda_{bc},\qquad
\tau=\frac{n_{abc}}{\Delta}.
\]
The positivity proved above gives $x,y,z>0$, and $\tau\ge0$.  The degree
equations give the common denominator identity
\[
1+\lambda_{ab}+\rho_c
=1+\lambda_{ac}+\rho_b
=1+\lambda_{bc}+\rho_a
=x+y+z+\tau.
\]
For example,
\[
1+\lambda_{ab}+\rho_c
=3-\frac{n_{ab}+n_{ac}+n_{bc}+2n_{abc}}{\Delta},
\]
and the two cyclic expressions are identical.  Thus, with
$M=x+y+z+\tau$,
\[
C_{ab}+C_{ac}+C_{bc}
=\frac2M\left(\frac1{1+x}+\frac1{1+y}+\frac1{1+z}\right).
\]
Since $\tau\ge0$, replacing $M$ by $x+y+z$ can only increase the expression.

For positive $x,y,z$,
\[
\begin{aligned}
&(x+y+z)\left(\frac1{x(1+x)}+\frac1{y(1+y)}
  +\frac1{z(1+z)}\right)
-3\left(\frac1{1+x}+\frac1{1+y}+\frac1{1+z}\right)\\
&=
\frac{(x-y)^2(x+y+1)}{xy(1+x)(1+y)}
+\frac{(x-z)^2(x+z+1)}{xz(1+x)(1+z)}
+\frac{(y-z)^2(y+z+1)}{yz(1+y)(1+z)} .
\end{aligned}
\]
Each term on the right is nonnegative, so
\[
\frac2{x+y+z}\left(\frac1{1+x}+\frac1{1+y}+\frac1{1+z}\right)
\le
\frac23\left(\frac1{x(1+x)}+\frac1{y(1+y)}+\frac1{z(1+z)}\right).
\]
Substituting back $x=1-\ell_{ab}$, $y=1-\ell_{ac}$, and
$z=1-\ell_{bc}$ yields
\[
C_{ab}+C_{ac}+C_{bc}
\le
\frac23\sum_{\{p,q\}\subseteq Q}
\frac1{(1-\ell_{\{p,q\}})(2-\ell_{\{p,q\}})}.
\]
Because
\[
\frac23\frac1{(1-\ell)(2-\ell)}
=\frac13\left(\frac{2}{(2-\ell)(1-\ell)}-1\right)+\frac13,
\]
and there are three pairs in $Q$, the last bound is exactly
\[
C_{ab}+C_{ac}+C_{bc}
\le
1+\frac13\sum_{\{p,q\}\subseteq Q}
\left(\frac{2}{(2-\ell_{\{p,q\}})(1-\ell_{\{p,q\}})}-1\right).
\]
Multiplying by the activation factor $\Delta^{-3}$ proves the claimed
fixed-triple survival bound.
\end{proof}
